\begin{textsample}{POS Dim 2 – profiled\_gpt – Score 6.00 – profiled\_gpt/t000942\_gpt.txt}  \label{ex:f2_pos_013}
so what we ’d actually be giving them isn't just the learning content itself, it ’s this whole \textbf{sense} of a really warm, friendly relationship between the students and the lecturers.
That ’s incredibly marketable for universities, because it taps into that old Oxford–Cambridge tutorial ideal – you know, that image of a small group, close contact, people really being known and talked to \textbf{rather} \textbf{than} just being numbers in a lecture theatre.
If we do the case studies properly, that ’s exactly what we ’re showing : students and academics in conversation, working through things together, and it feels human and approachable.
It ’s not just “ here is a video ”, it ’s “ here is what it looks like to study here, to be part of this little intellectual community ”.
And the key is that we ’re not just handing over random bits of video ; we ’re offering a really clear, sustainable process for creating those case studies.
A proper structure.
So a university can see, step by step, how we go from \textbf{idea} to finished piece : how we choose the topic, how we involve the students, how we frame the interaction with the lecturer, how it ’s edited, how it ’s presented online.
Because it ’s clearly formatted, they can understand what they ’re buying, and they can see how it can be repeated across \textbf{different} modules and departments.
It \textbf{becomes} a model they can plug into their existing teaching \textbf{rather} \textbf{than} a one‑off gimmick.
And then, on the student side, they ’re not just passive consumers.
They can go online and preview these case studies, get a feel for the style and the people, and then there ’s the possibility of them actually taking part in them as well.
That ’s quite powerful : you ’re selling the \textbf{idea} that you might end up in one of these pieces, talking to a lecturer, being part of the story.
So the offer to the universities is : high‑quality content, plus this very attractive picture of close student–lecturer relationships, all wrapped up in a robust, repeatable format for making case studies that students can see and maybe join in with.

% matched lemmas: become, different, idea, rather, sense, than
\end{textsample}
